\chapter{Planning}
\label{chapter:planning}

The thesis runs at 100\% commitment, nominally eight to nine months, in blocks of one to three weeks against the official TU Delft milestones. Progress meetings are held every two weeks, with findings and questions prepared in advance, and documentation is prioritised over presentations. The plan is treated as a living document and is updated at each progress meeting.

{\small
\begin{longtable}{p{1.9cm}p{6.0cm}p{2.5cm}p{3.4cm}}
    \caption{Block planning. The detailed Gantt chart is maintained separately in \texttt{Planning.xlsx}.}
    \label{tab:planning} \\
        \toprule
        Block & Work & Dates & Milestone \\
        \midrule
        \endfirsthead
        \toprule
        Block & Work & Dates & Milestone \\
        \midrule
        \endhead
        Orientation & Preliminary concept and first research questions & 18 to 26 Jun 2026 & \\
        Orientation & Literature review: mechanism survey, classification and gap; research questions fixed & 22 Jun to 17 Jul 2026 & Questions fixed 16 Jul \\
        Orientation & Decide the thesis method; write the methodology and the block planning & 13 Jul to 28 Jul 2026 & \\
        Orientation & Finalise the research proposal; register in MyCase & to 27 Aug 2026 & \textbf{Research proposal review, 27 Aug 2026} \\
        -- & Holiday & 31 Jul to 21 Aug 2026 & \\
        Research 1 & Parametric geometry models of the carried architectures and the switch-point sweep (sub-question 2) & 27 Aug to 25 Sep 2026 & \\
        Research 1 & ESATAN familiarisation model, in parallel and off the critical path & 14 Sep to 9 Oct 2026 & \\
        Research 1 & Trade-off: criteria first, first-order sizing, AHP with graphical cross-check; one concept selected (sub-question 1) & 21 Sep to 16 Oct 2026 & \\
        Research 1 & CAD model of the selected concept, with the Python model specialised to it; results documented; thesis outline & 12 Oct to 30 Oct 2026 & \\
        Research 1 & Midterm review & 2 to 6 Nov 2026 & \textbf{Midterm go / no-go; breadboard and latch-warping scope decided} \\
        Research 2 & ESATAN thermal cases (VLEO \SI{300}{\kilo\meter} cycling), thermoelastic and modal FEM; latch-warping study and breadboard only if scoped in & 9 Nov 2026 to 8 Jan 2027 & \\
        Research 2 & RSS error budget; passive-versus-active assessment against the 3-DOF budget (sub-question 3); complete draft & 4 Jan to 5 Feb 2027 & \\
        Research 2 & Green light review & 8 to 12 Feb 2027 & \textbf{Green light; minimum 20 working days before defence} \\
        Dissemination & Finalise the report; data and models to repositories; defence preparation & 15 Feb to 12 Mar 2027 & \\
        Dissemination & Defence & 15 to 19 Mar 2027 & \textbf{Defence, the official completion date} \\
        \bottomrule
\end{longtable}
}

The critical path runs through the parametric models. The switch-point sweep feeds the trade-off, the trade-off feeds the single selected design, and only then can the thermal and structural analysis begin. The latch-warping study is kept off that path, which is scoped in at the midterm only if there is room. The two points of float are the September to October overlap between the sweep and the trade-off, and the January overlap between the analysis and the error budget.

The orientation phase closes approximately four weeks later than originally planned, because the proposal work extended across the holiday period of 31 July to 21 August 2026. Phase 1 begins on 27 August 2026, when the review was rescheduled to, so the midterm, green light and defence dates are unchanged.

\bigskip

\noindent \emph{The survey behind \Cref{chapter:sota} holds approximately 300 catalogued papers and 111 structured notes. The reference list below is limited to the works cited in this document.}
